%***********************************************************************
%* File            :    MAIN.tex
%*
%* HAUPTDOKUMENT
%*
%* Dieses Dokument bindet alle untergeordneten Dokumente
%* ein. Es wird beim Aufruf von pdflatex als Quelle angegeben.
%*
%* Autoren         :    Daniel Hering, Tobias Wartzek 
%*                      
%* Datum           :    05.12.2009
%***********************************************************************



\documentclass[12pt,a4paper,twoside,openright,index=totoc,listof=totoc,DIV15,BCOR=8.25mm,headinclude,footinclude=false,headsepline,notitlepage,numbers=noenddot,bibliography=totoc,headings=normal,parskip=full]{scrbook}

%************************************************************************
% Das Package muss leider hier eingebunden werden da sonst in den 
% Befehlen für hypersetup keinerlei Umlaute zur Verfügung stehen
%************************************************************************

\usepackage[utf8]{inputenc}            % erlaubt direkte Eingabe Deutscher
                                         % Sonderzeichen in .tex Dateien
                                         %(latin1=iso8859-1 fuer Unix Editoren 
                                         % oder Windows Editoren die auch latin1
                                         % beherrschen)
                                         
%************************************************************************
% Dateiinfos eintragen (u.a. von hypersetup und für die Titelseite verwendet)
%************************************************************************
\usepackage{pdfpages}
% Titel der Arbeit 
\newcommand*{\doclongtitle}{Vergleich von Methoden des Maschinellen Lernens zur Detektion von Vorhofflimmern in Kontaktlosen Signalen}

% Art der Arbeit (Diplomarbeit, Studienarbeit, Bachelorarbeit, etc.)
\newcommand*{\docsubject}{Bachelorarbeit}

% Name des Autors
\newcommand*{\docauthor}{Nik Büttner}

% E-Mail Adresse des Autors
\newcommand*{\docauthoremail}{nik.buettner@rwth-aachen.de}

% Schlüsselworte der Arbeit 
\newcommand*{\dockeywords}{MedIT Latex Studienarbeit Diplomarbeit Vorlage}

% Name des Betreuers
\newcommand*{\docsupervisor}{Dipl.-Ing. Mustermann}

\newcommand*{\docauthorphone}{+49 151 70092028}

% Bild auf der Titelseite
%\newcommand*{\doctitlepic}{images/beispiel1}
\newcommand*{\doctitlepic}{images/frontCoverPic}

% folgende Befehle nicht ändern!!
\newcommand*{\dochomepage}{www.medit.hia.rwth-aachen.de}
\newcommand*{\docshorttitle}{\doclongtitle}

%**********************************************
%* Pakete (HEADER)                           *
%**********************************************

\include{HEADER}


%*************************************************************
%*************************************************************
%*																													**
%* 					Beginn des Dokuments 						  							**
%*																													**
%*************************************************************
%*************************************************************

\begin{document}
\begin{sloppypar}                       % schöner Blocksatz trotz langer Worte



%*************************************************************
%																														 *
%										 ENDE ALLER DEFINITIONEN 								 *
%																														 *
%*************************************************************


\frontmatter    												%römische Nummerierung aktivieren
\include{titelseite}
%\include{titelseite_alt}
\cleardoubleemptypage


%**************
%* Danksagung *
%**************

\chapter{Danksagung}
\thispagestyle{empty}
<!-- start slipsum code -->
Your bones don't break, mine do. That's clear. Your cells react to bacteria and viruses differently than mine. You don't get sick, I do. That's also clear. But for some reason, you and I react the exact same way to water. We swallow it too fast, we choke. We get some in our lungs, we drown. However unreal it may seem, we are connected, you and I. We're on the same curve, just on opposite ends.

Do you see any Teletubbies in here? Do you see a slender plastic tag clipped to my shirt with my name printed on it? Do you see a little Asian child with a blank expression on his face sitting outside on a mechanical helicopter that shakes when you put quarters in it? No? Well, that's what you see at a toy store. And you must think you're in a toy store, because you're here shopping for an infant named Jeb.

You see? It's curious. Ted did figure it out - time travel. And when we get back, we gonna tell everyone. How it's possible, how it's done, what the dangers are. But then why fifty years in the future when the spacecraft encounters a black hole does the computer call it an 'unknown entry event'? Why don't they know? If they don't know, that means we never told anyone. And if we never told anyone it means we never made it back. Hence we die down here. Just as a matter of deductive logic.
<!-- end slipsum code -->
\cleardoubleemptypage

%************************************************
%* Erklärung (hier muss nichts geändert werden) *
%************************************************

\include{erklaerung}
\cleardoubleemptypage

%************
%* Abstract *
%************

\chapter{Zusammenfassung}

In diesem Dokument finden Sie eine Sammlung von wertvollen Tipps und Empfehlungen, die Ihnen beim Verfassen Ihrer schriftlichen Abschlussarbeit am Lehrstuhl für Medizinischen Informationstechnik (MEDIT) helfen sollen. Es geht dabei um verschiedene wichtige Aspekte, die Sie bei der Erstellung Ihrer Arbeit berücksichtigen sollten, um eine qualitativ hochwertige wissenschaftliche Arbeit zu gewährleisten.

Gleichzeitig ist dieses Dokument auch eine Vorlage, die direkt für die schriftliche Arbeit genutzt werden kann.

Das Dokument beginnt mit Hinweisen zum Ablauf der Arbeit, um Ihnen einen Überblick über die einzelnen Phasen der Erstellung einer Abschlussarbeit zu geben. Anschließend wird der Aufbau der Arbeit behandelt, wobei auf die verschiedenen Abschnitte wie Einleitung, Hauptteil und Schluss eingegangen wird. Auch die inhaltliche Gestaltung der Arbeit wird thematisiert, wobei es darum geht, wie Sie Ihre Forschungsfrage präzise formulieren, Ihre Methodik beschreiben und Ihre Ergebnisse darstellen und interpretieren. Nicht zuletzt geben wir einige Hinweise zu häufig verwendeten Tools.

Ein weiterer wichtiger Aspekt, der in dem Dokument behandelt wird, ist der Umgang mit Quellen. Hierbei geht es um die korrekte Zitierung und Anerkennung von fremden Gedanken und Ideen, um Plagiate zu vermeiden. Auch die Bedeutung einer sorgfältigen Recherche und Auswahl von Quellen wird hervorgehoben.

Insgesamt dient das Dokument als umfassender Leitfaden für Studierende und gibt wertvolle Tipps und Empfehlungen, um den Schreibprozess zu erleichtern und die Qualität der Arbeit zu verbessern.



\cleardoubleemptypage
\thispagestyle{empty}
\cleardoubleemptypage

%**********************
%* Inhaltsverzeichnis *
%**********************


\addcontentsline{toc}{chapter}{Inhaltsverzeichnis} % Eintrag im Inhaltsverzeichnis erstellen
\tableofcontents                        % Inhaltsverzeichnis anlegen
\cleardoubleemptypage



%*********************
%* Symbolverzeichnis *
%********************

\include{symbolverzeichnis}
\cleardoubleemptypage

\mainmatter							% Arabische Nummerierung, Beginn des Hauptteils

%**************
%* Einleitung *
%**************

\include{einleitung}

%*************
%* Hauptteil *
%*************
 
\chapter{Theoretische Grundlagen}
\section{Medizinische Grundlagen}
\subsection{Vorhofflimmern}
\subsection{Herzfrequenz und Herzratenvariabilität}
\section{Kontaktlose Messmodalitäten}
\subsection{Kapazitives EKG (cECG)}
\subsection{Photoplethysmographie (PPG)}
\subsection{Ballistokardiographie (BCG)}
\section{Signalqualität und Artefakte}
\section{Maschinelles Lernen zur Klassifikation}
\subsection{Klassifikatoren}
\subsection{Merkmalsbasierte Klassifikation und Fusion}
\subsection{Evaluationsmetriken und Kreuzvalidierung}
\section{Stand der Forschung}
\chapter{Methoden}
\section{Datengrundlage und Verarbeitungsüberblick}
\section{Signalvorverarbeitung}
\section{Schlag- und Peak-Detektion}
\section{Merkmalsextraktion}
\section{Signalqualitätsindizes}
\section{Klassifikation und multimodale Fusion}
\section{Evaluationsstrategie}

\chapter{Ergebnisse}
\section{Schlagdetektion und HR-Schätzung gegen GT}
\section{Modellvergleich}
\section{Beitrag der Modalitäten (Ablation)}
\section{Signalqualitätsindizes}
\chapter{Diskussion}
\section{Einordnung der Hauptergebnisse}
\section{Fenster- vs. Patientenebene und Leckage}
\section{BCG als negatives Ergebnis}
\section{Limitationen}
\section{Ausblick}
\chapter{Fazit und Ausblick}
Hier folgt dann Text...

%**********
%* Anhang *
%**********

\cleardoublepage
\appendix
\chapter{Verzeichnis der Hilfsmittel}\label{sec:AnhangHilfsmittel}
\section{Programme}


\begin{tabular}{|p{4cm}|p{12cm}|}
	\hline 
Textverarbeitung	&  PDF Latex , TeX Live (\url{https://tug.org/texlive}) version 2022\\ 
	\hline 
Grafik	&  Corel Draw V12\\ 
	\hline 
Literaturverwaltung	& JabRef 5.11 \\ 
	\hline 
	&  \\ 
	\hline 
	&  \\ 
	\hline 
	&  \\ 
	\hline 
\end{tabular} 



\section{KI Werkzeuge}\label{sec:AnhangHilfsmittel_KI}
\begin{tabular}{|p{4cm}|p{12cm}|}
	\hline 
Titelbild	& erzeugt unter Verwendung von  \url{https://www.bing.com/images/create} \\ 
\hline 
Bild \ref{fig:beispiel1} rechts & erzeugt unter Verwendung von  \url{https://www.bing.com/images/create} \\ 
	\hline 
Kapitel \ref{sect:KItools}	& \url{https://chat.mistral.ai/}: Text mit Vorschl"agen generiert und einige Textpassagen umformuliert  \\ 
	\hline 
	&  \\ 
	\hline 
	&  \\ 
	\hline 
\end{tabular} 



\chapter{Latex Nutzung}
\section{Latex DO's \& DON'Ts}
						
Zum richtigen Umgang mit \LaTeX wird die PDF-Datei "`l2tabu"' empfohlen. In ihr sind die wichtigsten Regeln im Umgang mit \LaTeX und dem Koma-Script zusammengefasst (\url{www.mast.queensu.ca/~andrew/LaTeX/latex-dos-and-donts.pdf}).




\section{Grafiken mit TikZ und PGFPlots}

Eine typische Aufgabe in Abschlussarbeiten ist die Erstellung von Grafiken. Eine Möglichkeit ist die Verwendung eines speziellen Zeichenprogramms, wie zum Beispiel Inkscape oder CorelDraw. Aber gerade bei einfachen Blockschaltbildern können andere Lösungen auch überzeugen.

Ein Beispiel ist das Programmpaket TikZ. Hier können in einer Beschreibungssprache solche Grafiken komponiert werden.
Die Verwendung ist nicht ganz einfach, aber Ruben Sander hat mit seinem Betreuer Alfred Hülkenberg eine schöne Seminararbeit verfasst, die die Verwendung gut erklärt \cite{Sander2021}.

Eine weitere praktische Erweiterung dazu ist PgfPlot, die TikZ um die Möglichkeit erweitert schöne Plots aus Messdaten zu erzeugen. Es gibt sogar Tools, dies direkt aus Matlab oder Python zu realisieren. Hier hat Julian Bigalke mit seinem Betreuer Alfred Hülkenberg gleichermaßen eine sehr informative Seminararbeit verfasst \cite{Bigalke2023}.

Beide Manuale werden in diesem Vorlagenpaket im Zip Archiv mitgeliefert. 

\cite{abcde}


%**********************************************************
% Literaturverzeichnis (hier muss nichts geändert werden) *
%**********************************************************

\bibliographystyle{alphadin}    
%\bibliographystyle{plain}% plain dinat
\bibliography{literatur}%


\end{sloppypar}
\end{document}
